\subsection{交换星代数}

引理 8. --- 设 X 与 Y 是度量空间，其中 X 是完备的。设 f 是从 X 到 Y 中的一个映射，使得

\[
d (f (x), f (y)) \geqslant d (x, y)
\]

对 X 中一切 x 与 y 成立，且使得 f 的图像在 X × Y 中是闭的。那么 f 是一个闭映射。

设 F 是 X 的一个闭子集，\((y_{n})_{n\in\mathbf{N}}\) 是 \(f(\mathrm{F})\) 中收敛于 \(y\in Y\) 的一个序列。对每个 \(n\in N\)，取 \(x_{n}\in F\) 使 \(f(x_{n})=y_{n}\)。假设蕴含序列 \((x_{n})_{n\in\mathbf{N}}\) 是 Cauchy 序列；设 \(x\in X\) 是它的极限；由于 F 是闭的，它是 F 的元素。此外，我们有 \((x_{n},f(x_{n}))\to(x,y)\) 在 \(X\times Y\) 中成立。由于 f 的图像是闭的，由此得 \(y=f(x)\) 属于 \(f(\mathrm{F})\)。

引理 9. --- 设 A 是一个交换星代数。A 的每个特征标都是 Hermite 的，且 Gelfand 变换是从 A 到 \(\mathcal{C}_{0}(\mathsf{X}(\mathsf{A}))\) 中的星代数态射。

只需证明第一个断言（I, p.~38 的命题 7 与 I, p.~98 的引理 3）。设 \(\chi\) 是 A 的一个特征标。对 A 的每个 Hermite 元 \(y\)，我们有 \(\chi(y) = \mathcal{G}_{\mathrm{A}}(y)(\chi) \in \mathrm{Sp}'(y) \subset \mathbf{R}\) (I, p.~106, 命题 4)。从而特征标 \(\chi\) 是 Hermite 的 (I, p.~95, 命题 1)。

定理 1. --- 设 A 是一个交换星代数。Gelfand 变换是从星代数 A 到星代数 \(\mathcal{C}_{0}(\mathsf{X}(\mathsf{A}))\)（由 \(\mathsf{X}(\mathsf{A})\) 上在无穷远处为 0 的连续函数组成）上的一个等距同构。

Gelfand 变换是从 A 到 \(\mathcal{C}_{0}(\mathsf{X}(\mathsf{A}))\) 中的一个对合代数态射（引理 9）。设 B 是它的像。它是 \(\mathcal{C}_{0}(\mathsf{X}(\mathsf{A}))\) 的一个对合子代数。按定义，B 的元素分离 \(\mathsf{X}(\mathsf{A})\) 的点。设 \(\chi \in \mathsf{X}(\mathsf{A})\)。存在 \(x \in \mathsf{A}\) 使 \(\chi(x) \neq 0\)，于是 \(f = \mathcal{G}(x)\) 是 B 中使 \(f(\chi) \neq 0\) 的一个元素。依 TG, X, p.~40, 推论 2，子代数 B 因此在 \(\mathcal{C}_{0}(\mathsf{X}(\mathsf{A}))\) 中稠密。

对 A 的每个 Hermite 元 \(y\)，我们有 \(\| \mathcal{G}(y)\| = \varrho (y) = \| y\|\)（I, p.~38 的命题 7 与 I, p.~104 的公式 (2)），从而对每个 \(x\in \mathrm{A}\) 有等式

\[
\| x \| ^ {2} = \| x ^ {*} x \| = \| \mathcal {G} (x ^ {*} x) \| = \| \overline {{\mathcal {G} (x)}} \cdot \mathcal {G} (x) \| = \| \mathcal {G} (x) \| ^ {2}.
\]

于是 G 是等距的，因此它的像 B 是闭的（引理 8）。我们得出结论 \(\mathrm{B} = \mathcal{C}_{0}(\mathrm{X}(\mathrm{A}))\)。

推论 1. --- 设 A 是一个星代数，x 是 A 的一个正规元。那么 \(\|x\| = \varrho(x)\)。

由于 A 的由 x 与 \(x^{*}\) 生成的星子代数是交换的，可设 A 是交换的。在这种情形，推论由定理 1、I, p.~102 的例 2 以及 I, p.~24 的定理 1 得出。

推论 2. --- 设 A 是一个交换星代数。

\begin{enumerate}
\def\labelenumi{\alph{enumi})}
\item
  存在局部紧拓扑空间 X，使 A 同构于星代数 \(\mathcal{C}_{0}(\mathrm{X})\)；
\item
  设 B 是一个交换星代数。映射 \(\pi \mapsto \mathsf{X}'(\pi)\) 是从 A 到 B 的星代数态射组成的集合到从 \(\mathsf{X}(\mathsf{B})\) 到 \(\mathsf{X}(\mathsf{A})\) 中的真偏映射组成的集合上的一个双射 (I, p.~33, 定义 1)。
\end{enumerate}

定理 1 建立了第一个断言，第二个断言由 I, p.~34 的命题 3 与 I, p.~102 的例 2 得出。

推论 3. --- 设 A 是一个幺交换星代数。

\begin{enumerate}
\def\labelenumi{\alph{enumi})}
\item
  存在紧拓扑空间 X，使 A 同构于星代数 \(\mathcal{C}(\mathrm{X})\)；
\item
  设 B 是一个交换幺星代数。映射 \(\pi \mapsto \mathsf{X}(\pi)\) 是从 A 到 B 的幺星代数态射组成的集合到从 \(\mathsf{X}(\mathsf{B})\) 到 \(\mathsf{X}(\mathsf{A})\) 中的连续映射组成的集合上的一个双射。
\end{enumerate}

这由前文及 I, p.~35 的命题 4 得出。

注记. --- *设 G 为以局部紧空间为对象、以真偏映射 (I, p.~33, 定义 1) 为态射的范畴，并设 S 为交换星代数的范畴，其态射是对合代数态射。考虑从 S 到反范畴 G° 的函子，它把交换星代数 A 对应到 A 的非零特征标组成的局部紧空间 X(A)，把交换星代数的态射 φ: A → B 对应到连续映射 X'(φ)。定理 1 与推论 2 表明这个函子是一个范畴等价，且该函子的一个拟逆是这样的函子：它把局部紧拓扑空间 X 对应到交换星代数 C₀(X)。

类似地，推论 3 表明紧空间范畴的反范畴等价于幺交换星代数的范畴。*

